\section{三类 AI frontier：生成、行动与物理控制}

底层资源会进入不同应用形态。课堂把重点分成 generative、agentic 与 physical AI，这个分类的教学价值在于：三者的 failure mode 和 evidence 完全不同。生成内容可以离线评测；agent 会改变 workflow state；physical system 直接影响人和设备，因此验证、权限和恢复要求逐级上升。

\subsection{从输出质量到行动安全}

上一节已经说明三类 AI 系统拥有不同的行动半径；本节把这个差异转成验证要求。读图时沿着 generative、agentic、physical 的顺序观察：系统从“生成候选”扩展到“修改数字状态”，再扩展到“控制物理世界”，每一步都会增加权限、影响、发现延迟和恢复成本，因此不能沿用同一套 benchmark。

\lecturefigure{07-three-ai-frontiers.png}{Generative、agentic 与 physical AI 需要逐级增强的证据和问责。}{本地字幕 16:10--20:20；概念重绘。}

\begin{knowledgebox}{读图：同一个 benchmark 不能验证三层系统}
Generative AI 可看 factuality、utility 与 diversity；agentic AI 还要看 tool correctness、state transition、retry 和 exception handling；physical AI 必须加入 sensor failure、control stability、safe stop、human override 与现场法规。模型得分只是系统证据的一部分。
\end{knowledgebox}

\begin{importantbox}{行动半径越大，验证成本越高}
可把风险粗略理解为
\[
R \propto P(failure)\times I(impact)\times D_{detect},
\]
其中 $P(failure)$ 是失败概率，$I(impact)$ 是影响，$D_{detect}$ 是发现和纠正延迟。Agent 获得更多权限时，即便模型错误率不变，影响和恢复延迟也可能放大，因此要增加权限最小化、simulation、approval gate 与 rollback。
\end{importantbox}

\subsection{医疗：搜索加速不等于治理消失}

课堂以药物配方和机器人心脏移植说明 AI 进入医疗。KFSHRC 官方材料确认 2024 年完成 fully robotic heart transplant；这里的“fully robotic”描述手术路径和 robotic system 使用，不应改写成没有医生负责的 autonomous surgery。类似地，AI 可以缩短 candidate generation 和 formulation search，clinical trial、regulatory review 与长期安全证据仍然存在。

\lecturefigure{08-healthcare-boundaries.png}{医疗 AI 可压缩研究搜索，但不能跳过临床证据与责任边界。}{本地字幕 16:40--18:20；KFSHRC 官方全机器人心脏移植资料。}

\begin{knowledgebox}{读图：两条时间线不要相加}
左侧是 research loop：生成候选、预测性质、安排实验；右侧是 care delivery：患者选择、手术计划、医生控制、术后监测。AI 可以改善两条链的部分步骤，但每条链有不同数据、伦理、监管和责任主体。
\end{knowledgebox}

\begin{warningbox}{“十年变两年”应理解为课堂案例}
字幕中的周期数字没有足够材料确认其起点、终点和疾病项目。讲义只保留可迁移结论：AI 可能缩短配方与实验搜索，但不能据此声称完整 drug development 或 approval 周期按相同比例压缩。
\end{warningbox}

\subsection{能源：Agent 必须进入闭环，而不是生成建议}

能源行业拥有 sensor、maintenance、geology 和 operational history，也有高安全成本。讲者提到 corrosion 与 drilling，真正的系统问题是：模型建议怎样进入 work order、由谁批准、结果怎样回流、异常怎样停机。没有 outcome capture 的 agent 只能生成分析，无法形成可持续的 operational learning。

\lecturefigure{09-energy-agent-loop.png}{垂直 Agent 需要数据、建议、人工决策和结果反馈构成闭环。}{本地字幕 18:20--20:20；概念重绘。}

\begin{knowledgebox}{读图：为什么 domain data 仍然不够}
历史数据可能反映旧设备、旧策略与 selection bias；传感器缺失会制造 false confidence；成功 operation 往往没有详细标签。Agent 因此需要 uncertainty、operator rationale、counterfactual 和 incident review，而不是把 proprietary data 直接丢给通用模型。
\end{knowledgebox}

\teachervoice{课堂把国家的既有行业能力视为 AI use case 起点：能源国家不必只训练“大而全”的模型，可以先在高价值、高数据密度、可反馈的 workflow 建立优势。这比“先有模型，再找问题”更接近系统工程。}

\subsection{本章小结}

三类 AI frontier 的共同底座是模型与算力，真正差异却在行动半径、数据闭环和责任结构。越接近公共服务与物理世界，越不能用单一 benchmark 代表系统成熟度。

